\documentclass[12pt, twoside]{article}
%\documentclass[12pt, twoside]{article}
\usepackage[letterpaper, margin=1in, headsep=0.2in]{geometry}
\setlength{\headheight}{0.6in}
%\usepackage[english]{babel}
\usepackage[utf8]{inputenc}
\usepackage{microtype}
\usepackage{amsmath}
\usepackage{amssymb}
%\usepackage{amsfonts}
\usepackage[nomessages]{fp} %\FPeval{\var-name}{2*sin(pi/6)}
\usepackage{siunitx} %units in math. eg 20\milli\meter
\usepackage{yhmath} % for arcs, overparenth command
\usepackage{tikz} %graphics
\usetikzlibrary{quotes, angles, arrows, arrows.meta}
\usepackage{pgfplots}
\usepgfplotslibrary{fillbetween}
\pgfplotsset{compat=1.16}
\usepackage{graphicx} %consider setting \graphicspath{{images/}}
\usepackage{parskip} %no paragraph indent
\usepackage{enumitem}
\usepackage{multicol}
\usepackage{venndiagram}

\usepackage{fancyhdr}
\pagestyle{fancy}
\fancyhf{}
\renewcommand{\headrulewidth}{0pt} % disable the underline of the header
\raggedbottom
\hfuzz=2mm %suppresses overfull box warnings

\usepackage{hyperref}

\fancyhead[LE]{\thepage}
\fancyhead[RO]{Markscheme}
\fancyhead[LO]{La Scuola d'Italia / Dr.\ Huson / IB Math 12 \\* 9 November 2026}

\begin{document}
\subsubsection*{3.4 Pre-Test: Trigonometry Review --- Markscheme}

\begin{enumerate}[itemsep=0.2cm]

%--- Problem 1: Circle radius 8 cm, angle AOB = 2.1 rad: minor arc length; chord AB by cosine rule; area and perimeter of the shaded segment [11 marks] ---
\item[1.] [Maximum mark: 11]

Circle centre $O$, $r = 8$, $A\hat{O}B = 2.1$ radians.

\textbf{(a)} $l = 8 \times 2.1$ \hfill \textbf{M1}

$= 16.8$ cm \hfill \textbf{A1}

\emph{Wrong method: the degree formula, $\frac{2.1}{360} \times 2\pi
\times 8 = 0.293$, earns M0A0.}

\textbf{(b)} $AB^{2} = 8^{2} + 8^{2} - 2 \times 8 \times 8 \cos 2.1$
\hfill \textbf{M1A1}

$AB = 13.9$ cm \quad $(13.8787\ldots)$ \hfill \textbf{A1}

\emph{M1 for substituting into the cosine rule, A1 for correct
substitution. Accept the right-triangle method
$AB = 2 \times 8\sin 1.05$ for M1A1. Wrong method: degree mode gives
$AB = 0.293$ (0.293198\ldots), earning M1A1A0. $AB^{2} = 192.620\ldots$
given as the answer earns M1A1A0.}

\textbf{(c)} area of sector $= \dfrac{1}{2} \times 8^{2} \times 2.1
= 67.2$ \hfill \textbf{A1}

area of triangle $AOB = \dfrac{1}{2} \times 8 \times 8 \sin 2.1
= 27.6226\ldots$ \hfill \textbf{A1}

shaded area $= 67.2 - 27.6226\ldots$ \hfill \textbf{M1}

$= 39.6\ \text{cm}^{2}$ \quad $(39.5773\ldots)$ \hfill \textbf{A1}

\emph{M1 for subtracting the triangle from the sector (their values).
Accept 39.6 from the rounded values $67.2 - 27.6$. Wrong methods: degree
mode for the triangle, $\frac{1}{2} \times 64 \sin 2.1^\circ = 1.17$,
gives $66.0$, earning A1A0M1A0; the sector area alone (67.2) earns
A1A0M0A0; $r$ for $r^{2}$ in the sector, $\frac{1}{2} \times 8 \times
2.1 = 8.4$, earns A0.}

\textbf{(d)} perimeter $= 16.8 + 13.8787\ldots$ \hfill \textbf{M1}

$= 30.7$ cm \quad $(30.6787\ldots)$ \hfill \textbf{A1}

\emph{M1 for arc plus chord. FT from (a) and (b). Wrong method: adding the
two radii, $16.8 + 13.88 + 16 = 46.7$, or the sector perimeter
$16.8 + 16 = 32.8$, earns M0A0.}

\newpage

%--- Problem 2: SSS garden 45, 60, 80 m: largest angle by cosine rule, area; separate triangle PQ = 14, QR = 10, angle P = 38 deg: ambiguous case, two values of angle R, difference of the two possible lengths PR [12 marks] ---
\item[2.] [Maximum mark: 12]

Garden: $AB = 45$, $BC = 60$, $AC = 80$.

\textbf{(a)} The largest angle is $A\hat{B}C$, opposite the longest side
$AC$.

$\cos A\hat{B}C = \dfrac{45^{2} + 60^{2} - 80^{2}}{2 \times 45 \times 60}$
\hfill \textbf{M1A1}

$A\hat{B}C = 98.3^\circ$ \quad $(98.2514\ldots)$ \hfill \textbf{A1}

\emph{M1 for substituting into the cosine rule for any angle, A1 for the
angle opposite $80$. $\cos A\hat{B}C = -0.143518\ldots$ ($-\frac{775}{5400}$).
Wrong methods: the angle opposite $45$, $33.8^\circ$, or opposite $60$,
$47.9^\circ$, earns M1A0A0; radian mode gives $1.71$ (1.71481\ldots):
accept it only with the unit ``radians'' stated, A1A0 otherwise.}

\textbf{(b)} area $= \dfrac{1}{2} \times 45 \times 60 \sin 98.2514\ldots^\circ$
\hfill \textbf{M1}

$= 1340\ \text{m}^{2}$ \quad $(1336.02\ldots)$ \hfill \textbf{A1}

\emph{Accept any correct pairing of two sides with their included angle.
Accept 1336. FT from (a). Wrong method: an angle that is not included
between the two sides, e.g.\ $\frac{1}{2} \times 45 \times 60 \sin
33.8^\circ$, earns M0A0.}

Triangle $PQR$: $PQ = 14$, $QR = 10$, $Q\hat{P}R = 38^\circ$.

\textbf{(c)} $\dfrac{\sin P\hat{R}Q}{14} = \dfrac{\sin 38^\circ}{10}$
\hfill \textbf{M1}

$\sin P\hat{R}Q = 0.861926\ldots$

$P\hat{R}Q = 59.5^\circ$ \quad $(59.5335\ldots)$, \quad
$P\hat{R}Q = 120^\circ$ \quad $(120.466\ldots)$ \hfill \textbf{A1A1}

\emph{Accept $120.5^\circ$. Wrong methods: only the acute angle earns
M1A1A0; sides paired with the wrong angles, $\sin P\hat{R}Q =
\frac{10\sin 38^\circ}{14}$, gives $26.1^\circ$, M0A0A0.}

\textbf{(d)} \textbf{Method 1 (sine rule).}

$P\hat{Q}R = 180 - 38 - 59.5335\ldots = 82.4665\ldots^\circ$ or
$P\hat{Q}R = 180 - 38 - 120.466\ldots = 21.5335\ldots^\circ$
\hfill \textbf{M1}

$PR = \dfrac{10\sin 82.4665\ldots^\circ}{\sin 38^\circ} = 16.1$
\quad $(16.1024\ldots)$ \hfill \textbf{A1}

$PR = \dfrac{10\sin 21.5335\ldots^\circ}{\sin 38^\circ} = 5.96$
\quad $(5.96181\ldots)$ \hfill \textbf{A1}

difference $= 16.1024\ldots - 5.96181\ldots = 10.1$ cm
\quad $(10.1406\ldots)$ \hfill \textbf{A1}

\emph{M1 for a third angle found from either value in (c). Method 2,
cosine rule: $10^{2} = 14^{2} + PR^{2} - 2 \times 14 \times PR\cos
38^\circ$ (M1), so $PR^{2} - 22.0643\ldots PR + 96 = 0$, $PR = 16.1$ or
$5.96$ (A1A1), difference $10.1$ (A1). Accept $16.1 - 5.96 = 10.14$. FT
from (c). Wrong methods: one length only earns M1A1A0A0; the sum $22.1$
earns the final A0.}

\newpage

%--- Problem 3: Unit circle: exact sin(4pi/3), tan(2pi/3); sin given in quadrant II, find cos and tan exactly; solve 2cos x + sqrt3 = 0 on [0, 2pi] exactly [9 marks] ---
\item[3.] [Maximum mark: 9]

\textbf{(a)(i)} $\sin\dfrac{4\pi}{3} = -\dfrac{\sqrt{3}}{2}$ \hfill \textbf{A1}

\textbf{(a)(ii)} $\tan\dfrac{2\pi}{3} = -\sqrt{3}$ \hfill \textbf{A1}

\emph{Accept $-\frac{3}{\sqrt{3}}$. Decimals from the GDC ($-0.866$,
$-1.73$) earn A0: exact values are asked for. Common errors, each A0: the
wrong sign; $\cos\frac{4\pi}{3} = -\frac{1}{2}$ given for (i);
$-\frac{1}{\sqrt{3}}$ (ratio inverted) for (ii).}

\textbf{(b)(i)} $\cos^{2}\theta = 1 - \left(\dfrac{2}{3}\right)^{2}
= \dfrac{5}{9}$ \hfill \textbf{M1}

$\cos\theta = \pm\dfrac{\sqrt{5}}{3}$ \hfill \textbf{A1}

$\cos\theta = -\dfrac{\sqrt{5}}{3}$, since $\theta$ is in the second
quadrant \hfill \textbf{A1}

\textbf{(b)(ii)} $\tan\theta = \dfrac{2/3}{-\sqrt{5}/3} = -\dfrac{2}{\sqrt{5}}$
\hfill \textbf{A1}

\emph{Accept a right-triangle method (sides 2, $\sqrt{5}$, 3) for the
M1A1. The final A1 in (i) is for the negative sign. Accept
$-\frac{2\sqrt{5}}{5}$ in (ii). Wrong method: the GDC route
$\cos\left(\arcsin\frac{2}{3}\right) = +\frac{\sqrt{5}}{3}$ gives the
first-quadrant angle, earning M1A1A0, and then $\tan\theta =
+\frac{2}{\sqrt{5}}$, A0 unless FT. FT (ii) from the candidate's
$\cos\theta$. Inverted ratio $\frac{\cos\theta}{\sin\theta} =
-\frac{\sqrt{5}}{2}$ earns A0. Decimals ($-0.745$, $-0.894$) earn A0.}

\textbf{(c)} $\cos x = -\dfrac{\sqrt{3}}{2}$ \hfill \textbf{M1}

$x = \dfrac{5\pi}{6}$, \ $x = \dfrac{7\pi}{6}$ \hfill \textbf{A1A1}

\emph{Wrong methods: the GDC gives only $\frac{5\pi}{6}$ (2.618), earning
M1A1A0; $x = \frac{\pi}{6}, \frac{11\pi}{6}$ (sign of $\cos x$ wrong) earns
M0A0A0; $\frac{5\pi}{6}$ and $-\frac{5\pi}{6}$ ($-\frac{5\pi}{6}$ outside the
interval) earns M1A1A0. Decimals ($2.62$, $3.67$) earn A1A0 at most: the
question asks for multiples of $\pi$. Extra values outside the interval
earn the final A0.}

\newpage

%--- Problem 4: Ferris wheel h(t) = 16 - 14cos(pi t/6): a and d from min and max; b from the period; height at t = 2.5; GDC time above 25 m in one turn; describe transformations from y = cos t (with reflection) [13 marks] ---
\item[4.] [Maximum mark: 13]

$h(t) = d - a\cos(bt)$; lowest $2$ m, highest $30$ m, period $12$ minutes.

\textbf{(a)} $a = \dfrac{30 - 2}{2} = 14$ \hfill \textbf{A1}

$d = \dfrac{30 + 2}{2} = 16$ \hfill \textbf{A1}

\emph{Wrong methods, each A0: $a = 28$ (not halved); $d = 14$.}

\textbf{(b)} $b = \dfrac{2\pi}{12}$ \hfill \textbf{M1}

$b = \dfrac{\pi}{6}$ \quad $(0.523598\ldots)$ \hfill \textbf{A1}

\emph{Accept $0.524$. Wrong methods: $b = 12$ (the period used as $b$), or
$b = 2\pi \times 12 = 75.4$, earns M0A0; $b = 30$ (degree measure) earns
M1A0.}

\textbf{(c)} $h(2.5) = 16 - 14\cos\left(\dfrac{\pi}{6} \times 2.5\right)$
\hfill \textbf{M1}

$= 12.4$ m \quad $(12.3765\ldots)$ \hfill \textbf{A1}

\emph{FT from (a), (b). $b = 30$ in degree mode is consistent: M1A1.
Wrong method: $b = \frac{\pi}{6}$ in degree mode gives $2.00$, M1A0.}

\textbf{(d)} $16 - 14\cos\left(\dfrac{\pi}{6}t\right) = 25$ \hfill \textbf{M1}

\textit{GDC: graph $y = h(t)$ and $y = 25$ for $0 \le t \le 12$; find both
intersections.}

$t = 4.33$ \ $(4.33350\ldots)$ \ and \ $t = 7.67$ \ $(7.66649\ldots)$
\hfill \textbf{A1}

length $= 7.66649\ldots - 4.33350\ldots = 3.33$ minutes
\quad $(3.33298\ldots)$ \hfill \textbf{A1}

\emph{M1 for $h(t) = 25$, $h(t) > 25$, or a sketch with the line
$h = 25$. The first A1 needs both values. Accept 3.34 from $7.67 - 4.33$.
Wrong methods: one intersection (4.33) as the answer, M1A0A0; the time
below 25 m, $8.67$, M1A1A0. FT from the candidate's $h$.}

\textbf{(e)} Reflection in the $t$-axis (horizontal axis) \hfill \textbf{A1}

Vertical stretch, scale factor $14$ \hfill \textbf{A1}

Horizontal stretch, scale factor $\dfrac{6}{\pi}$ \quad $(1.90985\ldots)$
\hfill \textbf{A1}

Vertical translation $16$ up \hfill \textbf{A1}

\emph{``Describe fully'' needs words (markscheme-spec \S19): one mark per
transformation, with its type, direction and size. Accept the scale
factor $1.91$ and ``translation by $\binom{0}{16}$''. Reject ``move'', and
``flip'' with no axis. Vertical stretch with scale factor $-14$ earns the
stretch mark, not the reflection mark. Equations with no words earn at
most 1 mark. Order: reflection and vertical stretch (either order) before
the vertical translation; deduct 1 mark for a wrong order. Wrong method:
horizontal scale factor $\frac{\pi}{6}$, A0. FT from (a) and (b).}

\end{enumerate}
\end{document}
